\section{Discussion}\label{sec:discussion}
% Ch.4 revision (2026-09-29): the per-frame scale table, the scale-vs-depth figure and the use-case table moved to
% Results (tab:scale, fig:scalevsdepth, tab:exp1); this chapter is revised next (Ch.5).

\subsection{The teacher and the camera}\label{sec:teacher}
The laser teacher is not measurably better than the LiDAR teacher (Section~\ref{sec:finetune}). Because \ftfaro{} and
\ftlidar{} also differ in the number of views, the honest claim is that a phone's depth sensor can replace the laser as
the teacher at no measurable cost on these rooms---not that the two are equal. Fourteen extra LiDAR-only rooms do not
help measurably either; \ftmain{} is kept as the main model because it was best on validation and is the only recipe
whose labels scale without a laser. Knowing the camera does not rescue metric depth: Depth Pro with the true focal length is
far worse than the pretrained DA-v2 Metric, and even with the scale removed per frame its AbsRel stays several times
higher; the claim is about this model on this low-resolution indoor input, not about intrinsics-conditioned models in
general.

\subsection{Why fine-tuning works---and why it is not enough}\label{sec:whyft}
\textbf{Fine-tuning learns the domain's multiplier.} The pretrained model places every room too far, by a factor of
1.18--1.45; \ftmain{} is within $\sim$20\,\% in both directions (Table~\ref{tab:scale}). \ftmain{} without calibration
(0.134) is on par with the pretrained model given a per-room scale fitted on ground truth (0.127): what a one-off
calibration provides has moved into the head's weights, learned from LiDAR in other rooms. With the encoder frozen,
shape barely changes (oracle 0.035 vs.\ 0.042).
\textbf{The per-frame swing remains.} Within one room, frames at p5 and p95 still need factors $\sim$1.5$\times$ apart,
as before fine-tuning; with the right scale on every frame \ftmain{}'s AbsRel would drop to 0.051. That information is not in
a single image and does not come from a better teacher; it must come from metric data at run time. Fine-tuning and
per-frame metric points are therefore complementary.

\subsection{Scale follows the view, not time}\label{sec:drift}\label{sec:metric}
The relative model shows why one calibration per scene fails (Section~\ref{sec:scale}). Within one scan the central
90\,\% of frames need scales 2.6--4.9$\times$ apart; the correlation of scale with time flips sign between scenes
($-0.40$ to $+0.30$), so this is \emph{not cumulative drift}, whereas the correlation with the frame's median depth is
negative in every scene ($-0.27$ to $-0.75$, Fig.~\ref{fig:scalevsdepth}). The worst frames are
downward shots of table-tops at 0.5--\SI{0.8}{m} and frame-filling blank walls or cabinets: with no size cue the model
guesses at ``room scale,'' and a table at \SI{0.5}{m} lands at 1.5--\SI{2}{m} as a phantom surface. The metric model
behaves the same way: right on average up to a per-room bias, wrong per frame (GT/pred from 0.43 to 1.35 within one
room), but with almost no shift term---one scalar per frame recovers 97\,\% of its gap to the oracle. In 3D the same
surface seen from several frames is placed at several distances and TSDF averages them into a thick slab.

Rooms that are hard for LiDAR are not those that are hard for monocular depth: the standard bedroom 47333774 is the
easiest for LiDAR (\rsLIDARb) but second-hardest for mono$+$oracle (\rsORACLEb), because a small room scanned close to
blank walls gives the model no cue. The mirrored bathroom is hard for every source whose scale is right (LiDAR
\rsLIDARd, oracle \rsORACLEd), so a bathroom-renovation product should expect worse than the averages.

\subsection{Implications for a product}\label{sec:business}
Table~\ref{tab:exp1} translates the results into product tiers. iPad LiDAR suffices for furniture placement but not
for renovation measurement. On a device without LiDAR, neither one-off calibration nor a raw metric model serves any
tier; \ftmain{} raises recall@\SI{10}{cm} from \rsMETRICRten{} to \rsFTLIDARALLRten{} with nothing at run time, a large step
that is still short of the overview tier. With per-frame sparse points the same device reaches \SI{\rsSPARSECM}{cm},
close to furniture placement. The order of investment the data suggest is therefore: \textbf{(1) per-frame scale from
VIO points, (2) a fine-tuned model to remove the domain bias, (3) model size last}---investing in the model first
returns $<$\SI{3}{cm} on a \SI{22}{cm} base. Strategically, an app with Pro-device users can collect its own teacher
during normal use (RGB, LiDAR depth and pose arrive together from ARKit), replacing a \$20k laser scanner; our data
show no accuracy case for laser capture for this purpose. At 4~fps DA-v2 Large takes $3.6\times$ the video length on
a laptop and DA-v2 Small $0.6\times$; fusion is $<3\,\%$ of the time.
